\documentclass[10pt,a4paper]{amsart}
\usepackage[margin=19mm]{geometry}
\usepackage{amsmath,amssymb,mathtools}
\usepackage[scr=rsfs,cal=euler]{mathalfa}
\usepackage{xfrac}
\usepackage[unicode,hidelinks,bookmarks=false]{hyperref}
\newcommand{\ie}{i.e.\ }
\newcommand{\cf}{cf.\ }
\newcommand{\resp}{respectively\ }
\newcommand{\Subsection}{\subsection}
\providecommand{\nobreakdash}{\nobreak\hbox{-}}
\setlength{\emergencystretch}{2em}
\multlinegap0pt
\usepackage[shortlabels]{enumitem}
\usepackage{xcolor}
\usepackage{amssymb}
\usepackage{tikz}
\usepackage{bm}
\usepackage{mathtools}
\newtheorem{prop}{Proposition}
\newtheorem{lem}{Lemma}
\usepackage{cleveref}
\crefname{prop}{Proposition}{Propositions}
\crefname{lem}{Lemma}{Lemmas}
\DeclareMathOperator{\ch}{ch}
\newcommand\scalemath[2]{\scalebox{#1}{\mbox{\ensuremath{\displaystyle #2}}}}
\DeclareMathOperator{\vol}{vol}
\title{On Rankin-Selberg integral structures and Euler systems for $\mathrm{GL}_2\times \mathrm{GL}_2$}
\author{Alexandros Groutides}
\date{Source: arXiv:2407.01377v3 (CC BY 4.0)}
\begin{document}
\maketitle

\noindent\emph{Source and licence.} On Rankin-Selberg integral structures and Euler systems for $\mathrm{GL}_2\times \mathrm{GL}_2$. Version arXiv:2407.01377v3 (CC BY 4.0), distributed by arXiv under the Creative Commons Attribution 4.0 International licence, http://creativecommons.org/licenses/by/4.0/, as recorded in the arXiv licence metadata for this exact version and shown on the abstract page https://arxiv.org/abs/2407.01377v3. Full licence notice: \texttt{licenses/arxiv-2407.01377-CC-BY-4.0.txt}.\par\smallskip

\noindent\textit{Editorial scope.} Counted unit: the article's lem lattices (source0.tex lines 1220--1226). The article's complete original proof of it is delivered with it (source0.tex lines 1227--1264, one \texttt{proof} environment of 38 lines). No mathematics is written, repaired or re-derived: the statement and the proof are the article's own lines, in the article's own notation, and no retained line is substituted or re-flowed.  Retained uncounted as the source-local context the proof needs: lines 936--940; lines 1186--1209.  Nothing was omitted from any retained span, so the package carries no omission record.  No retained line contains a citation, so the wrapper carries no bibliography and no bibliography entry.  No retained line contains a figure environment, an image-inclusion command or drawing code, so no image is delivered and the package declares no asset dependency.  Editorial formatting only, in addition: an AMS \texttt{amsart} preamble that restates the article's own declarations and macros used by the retained lines, namely the environments \texttt{prop}, \texttt{lem} on separate counters, together with the standard packages those lines need.  The retained environments print in this document as Proposition 1, Lemma 1 in order of appearance, whereas the article prints the same items with the numbering of its own full document.  The complete native source of the article as distributed in the arXiv e-print for this exact version is bundled unchanged as \texttt{sources/161638/source0.tex}.
\begin{align}\label{eq:111}
\scalemath{0.95}{C_c^{\infty,1}(P_p\backslash G_p/G_p^\circ,\mathbf{Z}[1/p])=\left\{f\in C_c^\infty(P_p\backslash G_p/G_p^\circ,\mathbf{Z}[1/p])\ |\ f\ \text{is valued in}\ \begin{cases}
    (p-1)\mathbf{Z}[1/p],\ \text{on}\ Z_p(T_p\times \{1\}) \\
    \mathbf{Z}[1/p],\ \text{everywhere else.}
\end{cases}\right\}}\end{align}

\begin{prop}\label{prop from inv to coinv}
    There is a Hecke equivariant isomorphism
    $$\Phi: C_c^\infty(P_p\backslash G_p/G_p^\circ)\overset{\simeq}{\longrightarrow} C_c^\infty(G_p/G_p^\circ)_{P_p,\delta_P},\ \ch(P_pgG_p^\circ)\mapsto\frac{1}{\vol_{P_p}\left(P_p\cap g G_p^\circ g^{-1}\right)}\ch( gG_p^\circ)$$
    where $\vol_{P_p}$ denotes the volume with respect to right or left normalised Haar measure and is independent of such choice.
    \begin{proof}
    Firstly, we show that the volume factor in question is independent of the choice of left or right normalised Haar measure. Indeed, 
    \begin{align*}
        \vol_{P_p}(P_p\cap g G_p^\circ g^{-1},d^Rx)&=\int_{P_p\cap gG_p^\circ g^{-1}} d^Rx\\
        &=\int_{P_p\cap g G_p^\circ g^{-1}} \delta_P(x)\ d^Lx=\vol_{P_p}(P_p\cap g G_p^\circ g^{-1},d^Lx)
    \end{align*}
    where the last equality follows from the fact that $\delta_P$ is trivial on $P_p\cap gG_p^\circ g^{-1}$.
        We now check that the map $\Phi$ is well-defined. To do this, it clearly suffices to show that $\Phi(\ch(P_p g G_p^\circ))$ agrees with $\Phi(\ch(P_p yg G_p^\circ))$ for $y\in P_p$. Indeed, $\ch(P_p yg G_p^\circ)$ gets mapped to 
        \begin{align}\label{eq:26}
            \frac{1}{\vol_{P_p}(P_p\cap yg G_p^\circ g^{-1}y^{-1},d^Rx)}\ch(yg G_p^\circ)&=\frac{1}{\vol_{P_p}(P_p\cap yg G_p^\circ g^{-1},d^Rx)}\ch(yg G_p^\circ)\\
           \nonumber &=\frac{\delta_P(y)}{\vol_{P_p}(P_p\cap yg G_p^\circ g^{-1},d^Rx)}\ch(gG_p^\circ).
        \end{align}
        The first equality follows from the fact that $d^Rx$ is right invariant, and the second equality follows from the coinvariant action.
        Now, $\vol_{P_p}(P_p\cap yg G_p^\circ g^{-1},d^Rx)$ is given by $\int_{P_p\cap yg G_p^\circ g^{-1}} d^Rx= \int_{P_p\cap yg G_p^\circ g^{-1}} \delta_P(x) d^Lx$ which is in turn equal to $\delta_P(y) \vol_{P_p}\left(P_p\cap g G_p^\circ g^{-1}\right)$. From \eqref{eq:26} we now see that $\Phi$ is well-defined. The map $\Phi$ is a priori a morphism of $\mathbf{C}$-vector spaces.  Showing directly that $\Phi$ is Hecke equivariant using its current definition appears to be very hard. We will instead construct an inverse morphism to $\Phi$ which can be readily seen to be Hecke equivariant. The result will then follow at once. There is a natural morphism 
        $$\Psi: C_c^\infty(G_p/G_p^\circ)\longrightarrow C_c^\infty(P_p\backslash G_p/G_p^\circ),\ \xi\mapsto \left(f_\xi:g\mapsto \int_{P_p}\xi(xg)\ d^Rx\right)$$
        which is Hecke equivariant by construction. Now, for $y\in P_p$ $$f_{y\cdot\xi}(g)=\int_{P_p}\xi(y^{-1}xg)\ d^Rx=\int_{P_p}\xi(y^{-1}xg)\delta_P(x)\ d^Lx=\delta_P(y)\int_{P_p}\xi(xg)\delta_P(x)\ d^Lx=\delta_P(y)f_\xi(g).$$
        Thus, the map $\Psi$ factors through $C_c^\infty(G_p/G_p^\circ)_{P_p,\delta_P}$. Finally, for $g\in G_p$, the function $f_{\ch(gG_p^\circ)}$ is supported on $P_pg G_p^\circ$ and 
        $f_{\ch(gG_p^\circ)}(g)=\vol_{P_p}\left(P_p\cap g G_p^\circ g^{-1}\right).$ Hence $\Psi(\ch(gG_p^\circ))=\vol_{P_p}\left(P_p\cap g G_p^\circ g^{-1}\right)\ch(P_pgG_p^\circ)$ which shows that $\Phi$ and $\Psi$ are mutually inverse as required.
    \end{proof}
\end{prop}

\begin{lem}\label{lem identification of lattices}
    The map $\Phi$ identifies the lattices 
    \begin{align*}C_c^\infty(P_p\backslash G_p/G_p^\circ,\mathbf{Z}[1/p])&\simeq_{\Phi} \mathcal{J}(G_p/G_p^\circ,\mathbf{Z}[1/p])_{P_p,\delta_P}\\
    C_c^{\infty,1}(P_p\backslash G_p/G_p^\circ,\mathbf{Z}[1/p])&\simeq_{\Phi} \mathcal{J}_1(G_p/G_p^\circ,\mathbf{Z}[1/p])_{P_p,\delta_P}
    \end{align*}
    where $C_c^{\infty,1}$ is defined in \eqref{eq:111}.
\end{lem}

\begin{proof}
    The first identification follows at once from the definitions and \Cref{prop from inv to coinv}. Thus we deal with the second one. It is straightforward to show that every double coset $P_pg G_p^\circ$ is equal to $P_p(z_1\left[\begin{smallmatrix}
        p^m & \\
        & 1
    \end{smallmatrix}\right],z_2)G_p^\circ$ with $m\in \mathbf{Z}$ and $ z_i\in Z_p$, or $P_p(z_1\left[\begin{smallmatrix}
        p^m & p^{-n}\\
        & 1
    \end{smallmatrix}\right],z_2)G_p^\circ$ with $n\in\mathbf{Z}_{\geq 1}, m\in\mathbf{Z}, m>-n $ and $z_i\in Z_p$. By definition of the two lattices and the map $\Phi$,  it suffices to show that 
    \begin{align*}
        \tfrac{1}{p-1}\cdot \vol_{P_p}\left(P_p\cap g G_p^\circ g^{-1}\right)&= \vol_{P_p}\left(P_p\cap g G_p^\circ[p] g^{-1}\right),\ g= \left(z_1\left[\begin{smallmatrix}
        p^m & \\
        & 1
    \end{smallmatrix}\right],z_2\right)\\
     \vol_{P_p}\left(P_p\cap g G_p^\circ g^{-1}\right)&=\vol_{P_p}\left(P_p\cap g G_p^\circ[p] g^{-1}\right),\ g=\left(z_1\left[\begin{smallmatrix}
        p^m & p^{-n}\\
        & 1
    \end{smallmatrix}\right],z_2\right).
    \end{align*}
    We firstly treat the first case, in which we have $P_p\cap gG_p^\circ g^{-1}=P_p^\circ\cap \left[\begin{smallmatrix}
        p^m & \\
        &1
    \end{smallmatrix}\right] H_p^\circ \left[\begin{smallmatrix}
        p^m & \\
        &1
    \end{smallmatrix}\right]^{-1}.$ A quick matrix calculation shows that this is nothing more than $P_p^\circ$ if $m\leq 0$ and $\left[\begin{smallmatrix}
        \mathbf{Z}_p^\times & p^m\mathbf{Z}_p\\
        & 1
    \end{smallmatrix}\right]$ if $m>0$. The former has volume $1$ and the latter has volume $\frac{1}{p^m}$. On the other hand, if $m\leq 0$ then $P_p\cap g G_p^\circ[p] g^{-1}=\left[\begin{smallmatrix}
        1+p\mathbf{Z}_p & \mathbf{Z}_p\\
        & 1
    \end{smallmatrix}\right]$ which clearly has volume $\frac{1}{p-1}$. If $m>0$ then $P_p\cap g G_p^\circ[p] g^{-1}=\left[\begin{smallmatrix}
        1+p\mathbf{Z}_p & p^m\mathbf{Z}_p \\
        &1
    \end{smallmatrix}\right]$ which has volume $\frac{1}{(p-1)p^m}$ as required. For the second case, note that $m+n$ is strictly greater than zero. Another quick matrix computation shows that in this case the two subgroups have identical volume which is given (up to a power of $p$) by the volume of $\left[\begin{smallmatrix}
        1+p^{n+m}\mathbf{Z}_p & \mathbf{Z}_p\\
        & 1
    \end{smallmatrix}\right]$.
\end{proof}

\end{document}
